% Generated draft: requires mathematical and source review.
\documentclass[11pt]{article}
\usepackage[margin=25mm]{geometry}
\usepackage{fontspec}
\setmainfont{Latin Modern Roman}
\usepackage{amsmath,amssymb,mathtools,unicode-math}
\setmathfont{Latin Modern Math}
\usepackage{enumitem,needspace,xurl,hyperref}
\setlength{\parindent}{0pt}
\setlength{\parskip}{5pt}
\setlength{\emergencystretch}{3em}
\newcommand{\N}{\mathbb{N}}
\newcommand{\Z}{\mathbb{Z}}
\newcommand{\Q}{\mathbb{Q}}
\newcommand{\R}{\mathbb{R}}
\newcommand{\C}{\mathbb{C}}
\newcommand{\F}{\mathbb{F}}
\newcommand{\A}{\mathbb{A}}
\newcommand{\PP}{\mathbb{P}}
\newcommand{\E}{\mathbb{E}}
\newcommand{\eps}{\varepsilon}
\newcommand{\dd}{\,\mathrm{d}}
\begin{document}
\section*{UnsolvedMath ID: 701193}
\textbf{Source problem code:} catalogue entry
\subsection{Definitions and mathematical statement}
For a group \(K\), define its centre by
\[
Z(K)=\{z\in K:zk=kz\text{ for every }k\in K\}.
\]
For a finite non-abelian group \(K\), its non-commuting graph \(\Gamma_K\) is the simple undirected graph with vertex set \(K\setminus Z(K)\), in which distinct vertices \(x,y\) are adjacent exactly when \(xy\ne yx\). Thus \(\Gamma_G\cong\Gamma_H\) means that there exists a bijection
\[
f:G\setminus Z(G)\longrightarrow H\setminus Z(H)
\]
such that, for every pair of distinct vertices \(x,y\),
\[
xy\ne yx\quad\Longleftrightarrow\quad f(x)f(y)\ne f(y)f(x).
\]
For a group \(K\), put \(Z_0(K)=\{1\}\) and define \(Z_{i+1}(K)\) as the inverse image of \(Z(K/Z_i(K))\) under the quotient map. The group \(K\) is nilpotent if \(Z_c(K)=K\) for some integer \(c\geq0\). Put \(K^{(0)}=K\) and
\[
K^{(i+1)}=\langle a^{-1}b^{-1}ab:a,b\in K^{(i)}\rangle.
\]
The group \(K\) is solvable if \(K^{(d)}=\{1\}\) for some integer \(d\geq0\).

The remaining questions of the supplied record are the following two universal assertions, each quantified over all finite non-abelian groups \(G,H\):
\begin{enumerate}
\item[(b)] Is it true that
\[
\bigl(\Gamma_G\cong\Gamma_H\ \text{and }H\text{ is nilpotent}\bigr)
\quad\Longrightarrow\quad G\text{ is nilpotent}?
\]
\item[(c)] Is it true that
\[
\bigl(\Gamma_G\cong\Gamma_H\ \text{and }H\text{ is solvable}\bigr)
\quad\Longrightarrow\quad G\text{ is solvable}?
\]
\end{enumerate}
No equality of group orders or inequality between centre orders is assumed in these questions.

For part (b), the supplied comment of 2009 reports an affirmative answer under the additional hypothesis \(|H|=|G|\). The supplied comments of 2023 report an affirmative answer under the additional hypothesis \(|Z(H)|\geq |Z(G)|\), citing a 2024 publication and describing the use of an earlier result with the further assumption that \(C_H(h)\) is abelian for every \(h\in H\setminus Z(H)\), where
\[
C_H(h)=\{u\in H:uh=hu\}.
\]
These are supplied descriptions of partial results, not independently verified assertions here.

The record also retains archived part (a): for every pair of finite non-abelian groups \(G,H\), does
\[
\bigl(\Gamma_G\cong\Gamma_H\ \text{and }H\text{ is simple}\bigr)
\quad\Longrightarrow\quad G\cong H
\]
hold? Here a simple group is a nontrivial group whose only normal subgroups are \(\{1\}\) and the whole group. The supplied archive labels part (a) affirmatively solved modulo the classification of finite simple groups, with separate references for sporadic simple groups, alternating groups, and simple groups of Lie type. Its comment of 2009 also reports an affirmative result for ``groups with disconnected prime graphs''; the supplied wording does not explicitly identify which group bears this additional condition. This historical qualification is retained without choosing an interpretation or adding it to parts (b) and (c).
\subsection{Short English statement}
For finite non-abelian groups, does an isomorphism of their non-commuting graphs transfer nilpotence, or solvability, from one group to the other? The supplied record labels the problem partially solved, reports additional hypotheses giving affirmative answers to part (b), and archives the simple-group recognition question (a) as solved modulo the classification of finite simple groups. These source assessments require human review.
\subsection{Sources}
\begin{itemize}
\item Local UnsolvedMath registry: supplied complete source record, permanent ID \(701193\), exact source code \texttt{catalogue entry}, title \emph{Kourovka Notebook Problem 16.1}; upstream revision \texttt{d371ee917ac9647fbb06cd6fc529deec6f3dda82}. The supplied status is \texttt{partially\_solved}. The status review is dated \texttt{2026-10-16}; this is supplied dated curation, not an independently verified literature assessment.
\item \emph{The Kourovka Notebook (Unsolved Problems in Group Theory)}, 16th issue (2006), Problem 16.1, proposed by A. Abdollahi, S. Akbari, and H. R. Maimani; supplied as updated in the 21st issue (2026), \texttt{arXiv:1401.0300v48}. The supplied collection description names E. I. Khukhro and V. D. Mazurov as editors.
\item A. Abdollahi, S. Akbari, H. R. Maimani, \emph{J. Algebra} \textbf{298} (2006), 468--492. Cited by the supplied comment of 2009 for part (b) under \(|H|=|G|\).
\item H. Shahverdi, \emph{J. Algebra} \textbf{642} (2024), 60--64. Cited by the supplied comments of 2023 for part (b) under \(|Z(H)|\geq |Z(G)|\). The comment year and publication year are reproduced as supplied.
\item V. Grazian, C. Monetta, \emph{J. Algebra} \textbf{633} (2023), 389--402. Cited in the supplied comments of 2023 as an earlier result with the additional assumption that every non-central element of \(H\) has an abelian centralizer.
\item L. Wang, W. Shi, \emph{Commun. Algebra} \textbf{36} (2008), 523--528. Cited by the supplied comment of 2009 on archived part (a) for groups with disconnected prime graphs.
\item Ch. Khan', G. Ch\`en', S. Go, \emph{Siberian Math. J.} \textbf{49}, no. 6 (2008), 1138--1146. Supplied archive reference for the sporadic simple group cases of part (a); author spellings are retained from the record.
\item A. Abdollahi, H. Shahverdi, \emph{J. Algebra} \textbf{357} (2012), 203--207. Supplied archive reference for the alternating group cases of part (a).
\item R. M. Solomon, A. J. Woldar, \emph{J. Group Theory} \textbf{16}, no. 6 (2013), 793--824. Supplied archive reference for the simple groups of Lie type cases of part (a). The archive qualifies its overall affirmative answer as modulo the classification of finite simple groups.
\end{itemize}
\end{document}
